\section{Related work}
\label{sec:related}

\paragraph{Selective prediction and failure detection.}
Selective classification studies how a predictor and rejection rule trade coverage
for risk \citep{geifman2017selective,selectivenet,feng}. Recent work has also
emphasized that failure-detection conclusions depend on evaluation choices and
aggregation \citep{jaeger,values}. We adopt AURC as a ranking functional, but make
its loss index explicit because segmentation admits several non-equivalent notions
of error.

\paragraph{Segmentation confidence.}
Pixelwise entropy and maximum probability are common single-pass uncertainty
signals; converting them into an image-level score requires an aggregation rule.
ValUES systematizes this design space for semantic segmentation \citep{values}.
SDC is the closest metric-aligned predecessor to our work: it uses soft Dice between
the probability map and deployed mask and derives a Dice-specific approximation
\citep{sdc}. Our contribution is complementary. We retain the deployed action,
replace the Dice-specific surrogate by a loss-indexed expected-risk construction,
and state the joint working posterior required to evaluate nonlinear mask losses.

\paragraph{Random sets and geometric losses.}
Representing a fuzzy membership function through random level sets has a long
history in random-set theory \citep{goodman,molchanov}. Shared-level constructions
also arise in probabilistic contour and excursion-set analysis
\citep{poethkow,bolin}. We use this construction as a transparent computational
posterior, not as an identification result. Boundary-aware segmentation objectives
and evaluation measures address behavior missed by overlap alone
\citep{karimi2020}; our experiments ask whether a boundary-indexed confidence
improves selective ranking under a correspondingly defined boundary risk.
